\scanpage{295}
\mspage{- 11 -}

n° 246

similitude amenant $D$ sur $D'$ s'appellera \unsure{le|la} \uline{diangle} de $D$ et $D'$ et se notera $(D,D')$. Le diangle déterminé par la similitude $S$ est d'ordre $2$ dans le groupe $G_{\unsure{o|0}}/H$. Pour qu'un élément $X=\xi.1_E+\handwritten{\eta} S$ de $G_o$ détermine un diangle d'ordre $2$, il faut et il suffit qu'on ait $2\,\xi\handwritten{\eta}=0$ ; donc le diangle déterminé par $S$ est le seul élément d'ordre $2$ dans $G_o/H$ si $k$ n'est pas de caractéristique $2$.

Si $X=\xi.1_E+\handwritten{\eta} S$ est un élément de $G_o$, $\xi,\handwritten{\eta}$ ne sont pas tous deux nuls; posons $\handwritten{\tau}(X)=\handwritten{\eta}/\xi$ si $\xi\ne0$, et $\handwritten{\tau}(X)=\handwritten{\infty}$ si $\xi=0$ ; autrement dit, on pose en tout cas $\handwritten{\tau}(X)=\handwritten{\eta}/\xi$, le second membre étant entendu au sens du \uline{corps projectif} $k_{\handwritten{\infty}}$ attaché à $k$ (Alg., Chap. II, 2e éd., App. III, n° 5)

Pour que deux éléments $X$, $X'$ de $G_o$ appartiennent à la même classe suivant $H$, il faut et il suffit que $\handwritten{\tau}(X)=\handwritten{\tau}(X')$ ; pour qu'un élément $\handwritten{\tau}$ de $k_{\handwritten{\infty}}$ soit image par $\handwritten{\tau}(X)$ d'un élément de $G_o$, il faut et il suffit qu'on ait $\delta\handwritten{\tau}^2\ne1$. Autrement dit, $\handwritten{\tau}(X)$ détermine une \uline{bijection} de $G_o/H$ sur le complémentaire dans $k_{\handwritten{\infty}}$ de l'ensemble des racines de $\delta\handwritten{\tau}^2=1$ (ensemble à deux éléments si $\delta$ est un carré dans $k$, vide dans le cas contraire). On a la formule :

\[
\handwritten{\tau}(XX')=\frac{\handwritten{\tau}(X)+\handwritten{\tau}(X')}{1+\delta\handwritten{\tau}(X)\handwritten{\tau}(X')}\tag*{(9),}
\]

dont la vérification est immédiate, d'abord dans le cas où $\handwritten{\tau}(X)$, $\handwritten{\tau}(X')$, $\handwritten{\tau}(XX')$ sont finis, puis dans tous les autres cas.

L'homomorphisme canonique de $G_o$ sur $G_o/H$ induit sur $\handwritten{\Gamma_o}$ un homomorphisme de $\handwritten{\Gamma_o}$ dans $G_o/H$ dont le noyau est $\handwritten{\Gamma_o\cap} H=\{\pm1_E\}$, groupe à deux éléments si $k$ n'est pas de caractéristique $2$. Mais dans ce cas on peut aussi définir un isomorphisme (non moins \deleted{\unsure{c}}khânonique) entre ces groupes :

PROPOSITION 2. - \uline{Pour tout} $X\handwritten{\in}G_o$, $\mathbb{N}(X)^{-1}X^2$ \uline{est une rotation} ;

% archive sheet number: 294

